% Part: first-order-logic
% Chapter: models-theories
% Section: expressing-props-of-structures

\documentclass[../../../include/open-logic-section]{subfiles}

\begin{document}

\olfileid{fol}{mat}{exs}

\olsection{التعبير عن خصائص ال\printtoken{P}{structure}}

\begin{explain}
من مطالب هذا البحث التعبير عن شروط الدوال والعلاقات، وبيان
متى تحقق الدوال والعلاقات في بنية تلك الشروط. فقد نريد تمييز
ال!!{structure}s الموافقة لما قصدناه بمعاني ال!!{predicate}s
من ال!!{structure}s المخالفة له. ولنا أن نعيّن
ال!!{structure}s المقصودة كما نشاء؛ كأن نشترط في تأويل
ال!!{predicate}~$\le$ أن يكون ترتيبًا، أو نقتصر على تأويلات~$\Lang
L$ التي يكون مجالها من المجموعات ويكون مدلول $\Obj \in$ فيها
علاقة «هو !!a{element} في». ولكن أيمكن أداء هذه الشروط
ب!!{sentence}s من اللغة نفسها؟ أيّ الشروط على
!!a{structure}~$\Struct M$ يمكن أن تؤديه
!!a{sentence}، أو مجموعة من ال!!{sentence}s، في لغة~$\Struct
M$؟ ليست الشروط كلها مما تؤديه اللغة. فلا توجد مثلًا جملة
من~$\Lang L_A$ لا تصدق في !!a{structure}~$\Struct M$ إلا حيث
$\Domain M = \Nat$؛ فلا سبيل إلى التعبير عن أن «المجال لا يحتوي
إلا على الأعداد الطبيعية». وقد يمكن، مع ذلك، التعبير عن «خواص
بنيوية» لل!!{structure}s. فأي خواص ال!!{structure}s تؤديها
ال!!{sentence}s؟ وبعبارة أخرى، أي فئات من ال!!{structure}s
يمكن أن نحدّها بأنها البنى التي تصدق فيها !!a{sentence}، أو
مجموعة من ال!!{sentence}s؟
\end{explain}

\begin{defn}[نموذج مجموعة]
لتكن $\Gamma$ مجموعة من ال!!{sentence}s في لغة~$\Lang L$. نقول إن
!!a{structure}~$\Struct M$ \emph{نموذج لـ}~$\Gamma$ إذا كان
$\Sat{M}{!A}$ لكل $!A \in \Gamma$.
\end{defn}

\begin{ex}
فالجملة $\lforall[x][x \le x]$ تصدق في~$\Struct M$ إذا وفقط إذا
كانت $\Assign{\le}{M}$ انعكاسية. وأما الجملة
$\lforall[x][\lforall[y][((x \le y \land y \le x) \lif x = y)]]$،
فتصدق في~$\Struct M$ إذا وفقط إذا كانت $\Assign{\le}{M}$ مضادة
للتناظر. وأما الجملة $\lforall[x][\lforall[y][\lforall[z][((x \le y \land y \le z)
      \lif x \le z)]]]$، فتصدق في~$\Struct M$ إذا وفقط إذا كانت
$\Assign{\le}{M}$ متعدية. فإذن نماذج المجموعة
\begin{align*}
\{\quad &\lforall[x][x \le x], \\
   & \lforall[x][\lforall[y][((x \le y \land y \le
    x) \lif x = y)]], \\
   &\lforall[x][\lforall[y][\lforall[z][((x \le y
      \land y \le z) \lif x \le z)]]] \quad \}
\end{align*}
هي البنى التي تكون فيها~$\Assign{\le}{M}$ انعكاسية ومضادة
للتناظر ومتعدية، دون غيرها؛ أي البنى التي يكون فيها ترتيبًا جزئيًا.
فلنا أن نجعل هذه الجمل بديهيات
\emph{لنظرية الرتبة الأولى للترتيبات الجزئية}.
\end{ex}

\end{document}
